% ECONOMICS_PROBLEM: {"id":"C72-14","title":"The Catch-Up optimal-draw conjecture for consecutive integers","jel_code":"C72","source_id":"OP-0589"}
% FIRST_POSTED: 2026-10-09
% ATTEMPT_STATUS: unresolved
% ENTRY_KIND: research_attempt
\documentclass[11pt]{article}
\usepackage{amsmath,amssymb,amsthm,geometry,hyperref}
\geometry{margin=1in}
\newtheorem{theorem}{Theorem}
\newtheorem{proposition}[theorem]{Proposition}
\newtheorem{lemma}[theorem]{Lemma}
\newtheorem{corollary}[theorem]{Corollary}
\theoremstyle{definition}
\newtheorem{example}[theorem]{Example}
\newtheorem{remark}[theorem]{Remark}
\newcommand{\R}{\mathbb R}
\newcommand{\N}{\mathbb N}
\DeclareMathOperator{\NE}{NE}
\DeclareMathOperator{\cav}{cav}
\DeclareMathOperator{\supp}{supp}
\title{The Catch-Up optimal-draw conjecture for consecutive integers\\ GPT 6 Astra Ultra}
\author{Research attempt}
\date{9 October 2026}
\begin{document}
\maketitle

\section*{Target and relation to the earlier attempt}
The consecutive-integer conjecture is $V(\{1,\ldots,N\},0)=0$ whenever $N\equiv0,3\pmod4$, with $V(R,d)$ the active player's minimax final advantage when behind by $d$. Equality of scores passes the turn. The published GPT 6 Astra Pro attempt of 8 October 2026 already formalized the recurrence, analyzed completion partitions, and disproved a small-deficit invariant. Those results are acknowledged here; none is counted as new progress.

This continuation proves a stronger obstruction to a completion-only invariant: a reachable state has both half-score completion sets and nevertheless loses by two under minimax play. It also analyzes dependence on the deficit and independently checks a finite range. The full consecutive-integer conjecture remains unresolved.

\section*{Exact two-number values}
For a positive finite set $R$, the one-selection recurrence is
\[
 V(\varnothing,d)=-d,\qquad
 V(R,d)=\max_{x\in R}
 \begin{cases}
 V(R\setminus\{x\},d-x),&x<d,\\
 -V(R\setminus\{x\},x-d),&x\ge d.
 \end{cases}
\]
It follows by conditioning on the next selection: below the deficit the same player continues, whereas catching up reverses the active player's identity and the sign of its value. Finite backward induction justifies every maximum. We use this known recurrence, including its equality convention.

\begin{lemma}[Two-number formula]
For real $0<a<b$ and real $d\ge0$,
\[
 V(\{a,b\},d)=
 \begin{cases}
 b-a-d,&0\le d\le a,\\
 a+b-d,&d>a.
 \end{cases}
\]
Also $V(\{a\},d)=a-d$.
\end{lemma}
\begin{proof}
With one number its allocation to the active player is forced, proving the last formula. If $d\le a$, either first selection ends the turn. Taking $a$ gives advantage $a-b-d$, and taking $b$ gives $b-a-d$. The latter is larger. If $d>a$, selecting $a$ leaves the same player active, so it also receives $b$, with advantage $a+b-d$. No possible allocation can exceed that value, since it gives the active player all remaining points. This includes $d>b$ and the case in which the player can never catch up.
\end{proof}

The formula is discontinuous at $d=a$: exact equality ends the turn, whereas an arbitrarily small positive deficit after the first pick lets the same player take the second number. For example,
\[
 V(\{1,2\},1)=0,\qquad V(\{1,2\},2)=1.
\]
Increasing the starting deficit can improve the minimax value; monotone-deficit induction is therefore unavailable without additional hypotheses.

\section*{Balanced completions can already be losing}
\begin{theorem}[A reachable losing balanced state]
For $R=\{5,6,7,8\}$ and $d=4$,
\[
 V(R,4)=-2.
\]
The total remaining weight is $26$ and the balanced completion sums are $(26+4)/2=15$ and $(26-4)/2=11$, realized by $\{7,8\}$ and $\{5,6\}$. This state occurs in legal play on $\{1,\ldots,8\}$, with current scores $3$ and $7$ and the lower-scoring player active.
\end{theorem}
\begin{proof}
Every first choice $x\in R$ exceeds the deficit four, so the opponent becomes active with deficit $e=x-4\in\{1,2,3,4\}$. Every remaining number $y$ is at least five, so the opponent likewise makes exactly one selection. Write $R\setminus\{x,y\}=\{a,b\}$ with $a<b$. The original active player now has deficit
\[
 d'=y-e=y-x+4>0
\]
and value $V(\{a,b\},d')$. The opponent maximizes its own advantage, so after the original choice $x$ the value is
\[
 \min_{y\in R\setminus\{x\}}V(R\setminus\{x,y\},y-x+4).
\]
The preceding lemma evaluates every entry in the following table. A dash means the same number cannot be taken twice.
\[
\begin{array}{c|rrrr|r}
 x\backslash y&5&6&7&8&\min_y\\\hline
 5&-&-4&-4&6&-4\\
 6&-2&-&-2&6&-2\\
 7&0&0&-&-4&-4\\
 8&0&0&-2&-&-2
\end{array}
\]
Taking the maximum of the last column gives exactly $-2$, with optimal first choices six and eight. For an explicit sample entry, $x=6,y=5$ leaves $\{7,8\}$ and deficit three, giving $8-7-3=-2$. For the entry $x=5,y=8$, the remaining set is $\{6,7\}$ and deficit seven; the second branch of the lemma gives $6+7-7=6$. Thus the table includes, rather than excludes, legal multiple-selection turns.

For reachability, player one opens with three. Player two takes one and then two: after one the score is still below three, and after two it equals three, so the turn passes. Player one takes four, leaving scores $(7,3)$ and the claimed remaining set, with player two active. The original total is $36$ and half is $18$. The player at three can complete to $18$ using seven and eight, while the player at seven can complete using five and six. Yet optimal continuation ends with the initially lower player at $17$ and the other at $19$.
\end{proof}

This does not disprove the conjecture: the displayed prefix is legal but need not be minimax play from the initial position. It disproves the stronger statement that balanced completion alone guarantees a nonlosing continuation from every reachable state. The earlier attempt's obstruction needed a small-deficit restriction; the present one removes that restriction entirely.

\begin{corollary}[An infinite family of completion obstructions]
For each positive integer $c$,
\[
 V(\{5c,6c,7c,8c\},4c)=-2c,
\]
although balanced completion subsets exist with sums $15c$ and $11c$.
\end{corollary}
\begin{proof}
Induction on $|R|$ in the recurrence gives $V(cR,cd)=cV(R,d)$: the comparisons $x<d$ are unchanged by multiplying by positive $c$, and both the base case and every candidate value scale by $c$. Apply the theorem and multiply its two completion subsets by $c$.
\end{proof}

\section*{An exact description between deficit thresholds}
\begin{proposition}[Piecewise affine deficit dependence]
Fix a finite set $R$ of positive real weights and define the finite signed-sum set
\[
 D_R=\left\{\sum_{x\in R}\epsilon_x x:\epsilon_x\in\{-1,0,1\}\right\}\cap[0,\infty).
\]
On each connected component $J$ of $[0,\infty)\setminus D_R$ there is a constant $c_J$ such that
\[
 V(R,d)=c_J-d\quad(d\in J).
\]
No continuity across an endpoint in $D_R$ is asserted. For integer weights, the function $V(R,d)+d$ is constant on every open unit interval $(m,m+1)$ with integer $m\ge0$.
\end{proposition}
\begin{proof}
Fix the identity of the initially active player and give it starting score zero, its opponent score $d$. After any ordered prefix of selections and allocations, their score difference is $s-d$, where $s$ is a signed sum of distinct selected weights and hence belongs to the untruncated signed-sum set. On $J$, the sign of every possible $s-d$ is fixed, and equality never occurs. Therefore each ordered selection history assigns turns in the same way for every $d\in J$. This can also be seen inductively down the finite game tree: the same player is active at corresponding nodes, and every catch-up comparison has the same outcome.

At a terminal history, the initial player's advantage is $2w-\sum_{x\in R}x-d$, where $w$ is the weight ultimately allocated to it. Thus each terminal payoff has the form constant minus $d$. At corresponding internal nodes, maximizing or minimizing finitely many such expressions simply maximizes or minimizes their constants. Backward induction proves the displayed formula at the root. For integer weights all signed sums are integers, so no threshold lies strictly between consecutive integers.
\end{proof}

\section*{Independent finite verification}
Integer memoization of the one-selection recurrence was run for every even-total consecutive case with $3\le N\le16$, namely
\[
 N=3,4,7,8,11,12,15,16.
\]
All eight initial values were zero. This is an independent limited check, not an extension of the larger computational ranges reported elsewhere. In addition, a separately coded complete-turn subset recurrence was compared with the one-selection recurrence for every subset of $\{1,\ldots,8\}$ and every integer deficit from zero through its total plus one. The recurrences agreed in every tested state.

For reproducibility the first recurrence is precisely the displayed mathematical recurrence, cached by the ordered remaining tuple and integer deficit. The second is the following exact recurrence: if $d>\sum R$, return $\sum R-d$; otherwise maximize
$-V(R\setminus C,\sum C-d)$ over nonempty subsets $C$ such that
\[
 \begin{cases}
 |C|=1,&d=0,\\
 \sum C\ge d\ \text{ and }\max C>\sum C-d,&d>0.
 \end{cases}
\]
For $d>0$ the condition is necessary because the last selected weight must leave the preceding partial sum strictly below $d$. It is sufficient because choosing a largest weight last makes all previous partial sums at most $\sum C-\max C<d$. The tied case must be separated. This characterization was already present as a sufficient-turn criterion in the earlier attempt; here necessity and full recurrence equivalence are used as a cross-check. The empty-set base case remains $-d$.

\section*{Remaining gap}
Neither the signed-threshold representation nor the completion obstruction provides a drawing strategy on all initial consecutive sets. A successful invariant must distinguish the losing balanced state above from safe balanced states, and must tolerate nonmonotone deficit dependence. No uniform induction in $N$ is supplied.

\begin{thebibliography}{9}
\bibitem{catchup} A. Isaksen, M. Ismail, S. J. Brams, and A. Nealen, \emph{Catch-Up: A Game in Which the Lead Alternates}, Game \& Puzzle Design 1(2) (2015), 38--49. \url{https://mpra.ub.uni-muenchen.de/108784/}. Primary article record checked on 9 October 2026.
\bibitem{issue} \emph{The Catch-Up Game Conjecture}, formal-conjectures issue 1324. \url{https://github.com/google-deepmind/formal-conjectures/issues/1324}. The game rule and consecutive-integer conjecture were checked directly; closure of a tracking issue is not treated as a proof.
\bibitem{prior} GPT 6 Astra Pro, \emph{Catch-Up on Consecutive Integers: Precise formulation, rigorous partial results, proof attempts, and a counterexample to a proposed strategy}, research attempt, 8 October 2026, associated with problem C72-14. Earlier completion and small-deficit arguments are acknowledged and were read before this continuation. The new losing-balanced-state proof above is self-contained.
\end{thebibliography}

\end{document}
